\ExerciseSection

\begin{enumerate}
\def\labelenumi{\arabic{enumi})}
\tightlist
\item
  * a) 设 \(f\) 是加法群 \(\mathbf{R}\) 到其自身的同态。证明：若 \(f\) 的图象在 \(\mathbf{R}^2\) 中不稠密，则 \(f\) 形如 \(x \to ax\) 。{[}考察 \(\mathbf{R}^2\) 中 \(f\) 的图象的闭包，并利用 \(\mathbf{R}^2\) 的闭子群结构定理（第 VII 章，§ 1，第 2 小节，定理 2）{]}。 * {[}试与第 VI 章，§ 1，习题 12 b) 以及第 IV 章，§ 6，习题 2 比较。{]}
\end{enumerate}

\begin{enumerate}
\def\labelenumi{\alph{enumi})}
\setcounter{enumi}{1}
\tightlist
\item
  若 f 的图象在 \(R^{2}\) 中稠密，则 \(R^{2}\) 的拓扑在映射 \(x \to (x, f(x))\) 下的逆像与 R 的群结构相容，且严格细于 R 的通常拓扑。若此外 f 还是单射，则 R 的通常拓扑在 f 下的逆像与 R 的群结构相容，且与 R 的通常拓扑不可比较。
\end{enumerate}

\begin{enumerate}
\def\labelenumi{\arabic{enumi})}
\setcounter{enumi}{1}
\tightlist
\item
  设 \(\mathfrak{C}\) 是 \(\mathbf{R}\) 上的豪斯多夫拓扑，它与 \(\mathbf{R}\) 的群结构相容，且严格粗于通常拓扑 \(\mathfrak{C}_0\) 。
\end{enumerate}

\begin{enumerate}
\def\labelenumi{\alph{enumi})}
\item
  证明：关于 \(\mathfrak{C}\)，o 的每个开邻域在 \(\mathbf{R}\) 中都无界（注意，粗于紧空间拓扑的豪斯多夫拓扑必与该拓扑一致）。
\item
  设 \(V \neq R\) 是关于 \(\mathfrak{C}\) 的 o 的一个对称开邻域，并设 W 是关于 \(\mathfrak{C}\) 的 o 的一个对称开邻域，满足 \(W + W \subset V\) 。证明：若 a 是 V 的包含 o 的那个连通分支（关于 \(\mathfrak{C}_{0}\)）的长度，则 W 的每个连通分支（关于 \(\mathfrak{C}_{0}\)）的长度都 \(\leqslant a\) 。此外，R - W 的内部（关于 \(\mathfrak{C}_{0}\)）的各连通分支的长度所成的集是有界的{[}利用 a) 以及存在 o 的对称开邻域 \(W_{1}\) （关于 \(\mathfrak{C}\)）使 \(W_{1} + W_{1} \subset W\) 这一事实{]}。
\item
  由 b) 推出：\(\mathbf{R}\) 在拓扑 \(\mathfrak{C}\) 下是预紧的且不是局部紧的。
\item
  同样证明：\(\mathbf{Z}\) 在任何与其群结构相容且异于离散拓扑的拓扑 \(\mathfrak{C}\) 下都是预紧的且不是局部紧的。
\item
  设 \(f\) 是 \(\Gamma\) （\(\mathbf{R}\) 的子群，不只由 o 组成，赋有由 \(\mathfrak{C}_0\) 诱导的拓扑）到完备豪斯多夫群 G 中的连续同态。证明：若 \(f\) 不是 \(\Gamma\) 到 \(f(\Gamma)\) （\(\mathbf{G}\) 的子群）上的同构，则 \(f(\Gamma)\) 在 \(\mathbf{G}\) 中是相对紧的{[}利用 \(c\) ) 与 \(d)\){]}。
\end{enumerate}

* \(f\) ) 对每个整数 \(n \geqslant 2\) ，给出连续单射同态 \(f\) （\(\mathbf{R}\) 到 \(\mathbf{T}^n\) 中的）的例子，使得 \(f(\mathbf{R})\) 在 \(\mathbf{T}^n\) 中稠密（参见第 VII 章，§ 1，第 4 小节，命题 7 的推论 1）。 *

\begin{enumerate}
\def\labelenumi{\arabic{enumi})}
\setcounter{enumi}{2}
\tightlist
\item
  证明群 \(\mathbf{T}\) 代数同构于乘积 \(\mathbf{R} \times (\mathbf{Q}/\mathbf{Z})\)（在 \(\mathbf{R}\) 中取一个适当的哈梅尔基）。由此推出：存在 \(\mathbf{R}\) 上的一个豪斯多夫拓扑，它与群结构相容，与通常拓扑不可比较，并且 \(\mathbf{R}\) 关于它是预紧的。
\end{enumerate}

